% Generated draft: requires mathematical and source review.
\documentclass[11pt]{article}
\usepackage[margin=25mm]{geometry}
\usepackage{fontspec}
\setmainfont{Latin Modern Roman}
\usepackage{amsmath,amssymb,mathtools,unicode-math}
\setmathfont{Latin Modern Math}
\usepackage{enumitem,needspace,xurl,hyperref}
\setlength{\parindent}{0pt}
\setlength{\parskip}{5pt}
\setlength{\emergencystretch}{3em}
\newcommand{\N}{\mathbb{N}}
\newcommand{\Z}{\mathbb{Z}}
\newcommand{\Q}{\mathbb{Q}}
\newcommand{\R}{\mathbb{R}}
\newcommand{\C}{\mathbb{C}}
\newcommand{\F}{\mathbb{F}}
\newcommand{\A}{\mathbb{A}}
\newcommand{\PP}{\mathbb{P}}
\newcommand{\E}{\mathbb{E}}
\newcommand{\eps}{\varepsilon}
\newcommand{\dd}{\,\mathrm{d}}
\begin{document}
\section*{UnsolvedMath ID: 30000539}
\textbf{Source problem code:} catalogue entry
\subsection{Definitions and mathematical statement}
Let \(M\) be the closed oriented surface in the source and \(g\) an orientation-preserving mapping class. The genus, area form and precise Floer/monotonicity setup are not stated. The source's Conjecture 5 asks first for a monotone representative \(\phi\) with the printed decomposition
\[
 HF^*(\phi)\cong H^*(M_{\mathrm{id}},\partial M_{\mathrm{id}};\Z_2)
 \oplus\Z_2^{\,N(\phi|_{M\setminus M_{\mathrm{id}}})}.
\]
The equality is written as an isomorphism of the intended homology objects. Here \(HF^*\) is symplectic Floer cohomology with the displayed coefficients, and \(N(f)\) conventionally denotes the Nielsen number: the number of essential fixed-point classes of \(f\), with nonzero fixed-point index. The rightmost summand has one vector-space generator per counted class, but its grading is not supplied. The record does not specify whether the assertion is graded, or merely an isomorphism of total vector spaces.

The subspace \(M_{\mathrm{id}}\) is not defined in the record. It presumably refers to identity pieces in an intended surface decomposition, but the decomposition, boundary conventions, choice relative to \(g,\phi\) or a canonical representative, and invariance needed to define the restricted map are absent. It is not silently equated with the full pointwise fixed set of \(\phi\). Monotone also requires the intended symplectic mapping-torus/cohomological convention; no flux or genus hypothesis is invented. For the first assertion the quantifiers are: for every intended \(g\), there exists a suitable representative \(\phi\), not every representative has the displayed decomposition.

For a Thurston canonical representative \(\psi\), the second request asks about the printed identity
\[
 F^\infty(g)=\lambda=h(\psi)
 =\limsup_{n\to\infty}|N(\psi^n)|^{1/n},
\]
where \(\lambda\) is described as the largest pseudo-Anosov stretching factor. The original wording also describes \(F^\infty\) as asymptotic Floer-homology growth. A standard candidate is \(\limsup_n(\dim_{\Z_2}HF^*(\phi^n))^{1/n}\), when compatible iterates and the invariant are defined, but the record does not define \(F^\infty\) or any convention replacing a zero growth value by one. That candidate is not imposed as the missing definition.

There is an essential normalization issue in the printed chain: if \(h(\psi)\) is ordinary topological entropy in its logarithmic convention, a pseudo-Anosov stretch factor corresponds to \(\exp(h(\psi))\), while entropy corresponds to its logarithm. Thus the equality with \(h(\psi)\) cannot be silently interpreted as the standard entropy statement; either its notation means exponentiated entropy or a logarithm/exponential has been omitted. The displayed source identity is retained and this ambiguity is flagged rather than repaired. Components permuted by \(\psi\) require a per-iterate normalization of return-map stretch factors, and mapping classes with no pseudo-Anosov piece require a convention for \(\lambda\); neither is supplied. The limsup is not strengthened to an ordinary limit.

The literature triage checked 2026-08-21 reports the growth identity still conjectural in later literature, with finite-type and special mapping-class cases proved. Their exact hypotheses are absent. The open label and verification note reporting direct transcription of Conjecture 5 remain supplied metadata, not independently checked current status.
\subsection{Short English statement}
Seek a monotone representative with the stated Floer--Nielsen decomposition, and compare asymptotic Floer and Nielsen growth with pseudo-Anosov stretching. The full printed formulas are retained. The entropy equality needs a logarithm/exponential convention, while identity pieces, grading, genus and iteration conventions are missing.
\subsection{Sources}
\begin{itemize}
\item Local UnsolvedMath registry, permanent ID 30000539, source code catalogue entry, ``Floer Homology and Growth of Surface Mapping Classes,'' including the original growth description, both cleaned formulas and verification reporting Conjecture 5.
\item \emph{Geometric Group Theory, Hyperbolic Dynamics and Symplectic Geometry (2007)}, Oberwolfach Reports, pp.\ 1991--2058, DOI 10.4171/owr/2006/33: \url{https://doi.org/10.4171/owr/2006/33}. The citation uses 2007, while the DOI and proposed-year metadata use 2006.
\item Supplied reference ``Fel'shtyn, Nielsen theory and Floer homology'': \url{https://www.impan.pl/shop/en/publication/transaction/download/product/85845}.
\item Literature triage checked 2026-08-21; local record updated 2026-09-16. No external source has been visited or independently checked here.
\end{itemize}
\end{document}
